\NativeSlot{46}{43}{%
% p006_s01

\ExamQuestion{7.}{5}%
导航显示，从家到公园全程1800米，小明先以每分钟80米的速度走了5分钟，剩下的路程要在15分钟内走完，剩下的路程每分钟需要走多少米？综合算式：\BookBlank{}。\par
}
\NativeSlot{124}{73}{%
% p006_s02

\ExamQuestion{8.}{7}%
学校体育室准备给同学们发放运动器材。①一共有15盒羽毛球。②每盒12个。③已经分给同学们80个。④剩下的羽毛球要分给5个班级。平均每个班分几个？通过信息①和信息②可以求出\mbox{（\hspace{1.8em}）}，然后与信息③合并可以求出\mbox{（\hspace{1.8em}）}，最后再求出剩下的羽毛球平均每个班级能分到\mbox{（\hspace{1.8em}）}个。\par
}
\NativeSlot{251}{36}{%
% p006_s03

\ExamQuestion{9.}{5}%
\(\square+\square+\square+\square=100\)，\allowbreak{}\((\triangle+\triangle)\times\square=100\)，\allowbreak{}\(\square\times\triangle\times\triangle=\)\mbox{（\hspace{1.8em}）}。\par
}
\NativeSlot{277}{105}{%
\ExamGroup{二、选择（将正确答案的序号填在括号里）。{\normalfont\kaishu\fontsize{10.8bp}{14bp}\selectfont （每空3分，共21分）}}

% p006_s04

\ExamQuestion{1.}{5}%
给算式\(600\div24+6\times2\)添加括号，使运算顺序变为“先求和，再求积，最后求商”，下列添加方式正确的是\mbox{（\hspace{1.8em}）}。\par
\ExamOptions{A. \(600\div24+(6\times2)\)}{B. \(600\div(24+6)\times2\)}{C. \(600\div[(24+6)\times2]\)}{D. \((600\div24+6)\times2\)}
}
\NativeSlot{457}{57}{%
% p006_s05

\ExamQuestion{2.}{5}%
将“\(520+360=880\)”，“\(900-880=20\)”，“\(160\div20=8\)”三个算式合并成一个综合算式，正确的是\mbox{（\hspace{1.8em}）}。\par
\ExamOptions{A. \(900-(520+360)\div20\)}{B. \(160\div[900-(520+360)]\)}{C. \(160\div(900-520+360)\)}{D. \(160\div(900-520-360)\)}
}
\NativeSlot{561}{36}{%
% p006_s06

\ExamQuestion{3.}{5}%
\(25\times8\div25\times8=\)\mbox{（\hspace{1.8em}）}，\allowbreak{}\(25\times8\div(25\times8)=\)\mbox{（\hspace{1.8em}）}。\par
\ExamOptions{A. 1}{B. 64}{C. 0}{D. 625}
}
\NativeSlot{610}{36}{%
% p006_s07

\ExamQuestion{4.}{5}%
被减数、减数与差的和是120，那么被减数是\mbox{（\hspace{1.8em}）}。\par
\ExamOptions{A. 120}{B. 60}{C. 40}{D. 条件不够}
}
\NativeSlot{665}{77}{%
% p006_s08

\ExamQuestion{5.}{5}%
要使\(8\times(\square\div3)\div2=24\)，下列说法和计算结果正确的是\mbox{（\hspace{1.8em}）}。\par
\ExamOptions{A. 先算\(8\div2=4\)，再算\(4\times24=96\)，最后算\(96\div3=32\)}{B. 先算\(24\times2=48\)，再算\(48\div8=6\)，最后算\(6\times3=18\)}{C. 先算\(24\times2=48\)，再算\(48\div8=6\)，最后算\(6\div3=2\)}{D. 先算\(8\div2=4\)，再算\(24\div4=6\)，最后算\(6\div3=2\)}
}
